%% =====================================================================
%%  C-04 -- case interlude, typeset after P6
%%  -------------------------------------------------------------------
%%  A case is synthesis, not doctrine: it runs existing entries against
%%  one situation so their interaction is visible. Every move in the
%%  lattice names the entry doing the work. No new claims, no sources
%%  of its own (rule R23). Word band: 300-1000 (soft 850).
%%  =====================================================================
%% @kind: case
%% @id: C-04
%% @title: The Fake Job Offer
%% @strapline: Six signals, one payload --- and the two moves that would have cost nothing
%% @after-part: P6
%% @order: 10
%% @combines: T-18-01, T-15-01, T-10-01, T-23-03, T-24-01, T-24-05
%% @status: stable
%% @updated: 2026-09-19

\case{C-04}{The Fake Job Offer}{Six signals, one payload --- and the two moves that would have cost nothing}

\begin{story}
Marcus, a logistics analyst, gets a LinkedIn message from a recruiter at \emph{Meridian Freight Systems}. Everything checks at a glance: the logo, the sign-off block, her four thousand connections, a keynote photo of the managing director. \emph{Your colleague Daniel placed us}, she writes --- Daniel is in his network, though Daniel, asked later, has never heard of them. The role is pitched at a 38 percent raise, fully remote, and the process is already warm: \emph{final round Monday; we are choosing between two candidates}. The offer letter arrives Thursday --- signed in seventy-two hours --- and onboarding begins before any contract does. First, a refundable equipment deposit, nine hundred pounds, wired to account details that turn out to be a personal name. Second, the real payload: \emph{bring your current employer's client rate card so we can benchmark your book}. Marcus asks for the deposit terms in writing and a joint call with Daniel's placement manager. The terms arrive as a screenshot. The call slips twice, then goes silent. The Monday deadline passes; the two candidates were, as far as anyone can tell, always zero and one.
\end{story}

\begin{lattice}
  \casemove{Authority at template cost}{T-18-01} Logo, keynote, footer: the visible signs of an institution, produced for the price of a template. Deference fired on the symbols before any of them were checked.
  \casemove{A crowd of connections}{T-15-01} Four thousand connections and a colleague's name stood in for a reputation. The number did the convincing; the name did the vouching; neither had substance behind it.
  \casemove{A Monday that never existed}{T-10-01} The final-round deadline manufactured a loss frame around an offer of fixed scarcity --- theirs, not his. The clock was there to keep the deposit and the rate card from being examined in daylight.
  \casemove{One credential, tested}{T-23-03} The placement-manager call was the cheapest possible verification: one request, transcribable, refuseable. The lattice collapsed at exactly that node --- the credential could not survive being checked.
  \casemove{An invoice for goods never ordered}{T-24-01} The deposit and the rate card were debts created without consent --- unauthorised obligations arriving pre-addressed. He owed a stranger nothing, and the plainness of that sentence was the whole reply.
  \casemove{The standing no}{T-24-05} His refusal was a policy --- \emph{I don't pay to be paid} --- not a counter-negotiation. Policies have nothing to argue with; that is what makes them safe to hold.
\end{lattice}

\begin{lesson}
No single signal was conclusive --- real recruiters are often exactly this pushy, which is what the design borrows. The lattice is the finding: authority, crowd, clock and debt all aimed at one payload, the rate card, and every one of them flinched from the same two cheap tests --- one verification request and one standing no. Detection here was not a feeling or a read; it was two moves that cost an afternoon and declined to cost a career.
\end{lesson}
